\chapter{Manual for EREPOPT}
\label{chap:erepopt}

``\textbf{erepopt}'' is a program to optimize all DFTB parameters simultaneously.
To run the program, type: 

{\scriptsize erepopt erep.inp}

where \textbf{erepopt} is program name and \textbf{erep.inp} is name of the
input file. The input file is organized in ``block'' sections. Each section
begins with ``\textbf{\$blockname:}'' and end with ``\textbf{\$end:}''.
`\#' is the comment character. Everything after `\#' will be skipped.

\section{Overview}

\textbf{erepopt} couples a genetic algorithm (GA) optimizer with external DFTB
parameter generation and scoring tools. In each GA step, a candidate parameter
set is translated into DFTB Slater--Koster (\texttt{.skf}) files (via \texttt{skgen}
and optional post-processing), then evaluated with one of the scoring modes
selected by \textbf{fit\_type}:

\begin{description}
\item[\is{fit\_type = 0}] Fit repulsive potentials by calling \textbf{repopt} internally.
This is the default mode and requires a valid \textbf{rep.in} input for \textbf{repopt}.

\item[\is{fit\_type = 1}] Fit band energies / eigenvalues for each compound against a
reference list.

\item[\is{fit\_type = 2}] Fit density of states (DOS) curves against reference DOS data.
This mode runs \textbf{dftb+} and \textbf{dp\_dos}.
\end{description}

\textbf{Parallelization.} \textbf{erepopt} is compiled with MPI and typically
run with multiple ranks. Inside each rank it can also use OpenMP for external
tools via \textbf{nthreads} (exported as \texttt{OMP\_NUM\_THREADS}).

\section{Input}

\subsection{\$system:}
\begin{b4tableh}
   nthreads        & integer &$\geq 1$  &    1                          \\
   dftbversion     & string  &          &    dftb+                      \\
   skgen           & string  &          &    skgen                      \\
   onecent         & string  &          &    hfatom\_spin               \\
   twocent         & string  &          &    sktwocnt\_lr               \\
   skdefversion    & integer &          &    (no default)               \\
   xcfunctional    & string  &          &    (no default)               \\
   omega           & Real    &          &    (unset)                    \\
   superposition   & string  &          &    (no default)               \\
   repopt          & string  &          &    repopt                     \\
   power           & integer &$\geq 2$  &    2                          \\
   dgrid           & Real    &$\geq 0.0$&    0.1                        \\
   ngrid           & integer &$\geq 1$  &    100                        \\
   grids           & string  &          &    grids                      \\
   rep.in          & string  &          &    rep4e.in                   \\
   libdir          & string  &          &    libskf4e                   \\
   libadddir       & string  &          &    (unset)                    \\
   repadddir       & string  &          &    (unset)                    \\
   scratchfolder   & string  &          &    /tmp                       \\
   skfclean        & bool    &  0|1     &    1                          \\
   outfile         & string  &          &    erepopt.log                \\
   popinitialfile  & string  &          &    pop.initial.dat            \\
   popfinalfile    & string  &          &    pop.final.dat              \\
\end{b4tableh}

\begin{description}
\item[\is{nthreads}] Number of OpenMP threads used for external calls made by \textbf{erepopt}
(e.g. \texttt{skgen} and \textbf{repopt}). This value is exported as \texttt{OMP\_NUM\_THREADS}.
\item[\is{dftbversion}] DFTB executable used in \textbf{fit\_type = 1/2} (default: \texttt{dftb+}).
\item[\is{skgen}] The \texttt{skgen} executable.
\item[\is{onecent}] The one-center executable used by \texttt{skgen} (default: \texttt{hfatom\_spin}).
\item[\is{twocent}] The two-center executable used by \texttt{skgen} (default: \texttt{sktwocnt\_lr}).
\item[\is{skdefversion}] Version number written to \texttt{SkdefVersion} in the generated \texttt{skdef.hsd}.
\item[\is{xcfunctional}] XC functional string written to \texttt{XCFunctional} in \texttt{skdef.hsd} when \is{omega} is not used.
\item[\is{omega}] If provided, enables long-range corrected mode and writes
\texttt{func = xchybrid \{ omega = ... \}} and \texttt{LC omega} markers in \texttt{.skf} files.
\item[\is{superposition}] Value written to \texttt{Superposition} in \texttt{skdef.hsd}.
\item[\is{repopt}] Path to \textbf{repopt} executable (used in \textbf{fit\_type = 0}).
\item[\is{power}] Power used in \texttt{skdef.hsd} (\texttt{PowerCompression}). See \is{\$atomparameters:}.
\item[\is{dgrid}] \textbf{Deprecated / no effect.} Kept only for backward compatibility; the current implementation
reads this value but never uses it in any calculation or external call.
\item[\is{ngrid}] \textbf{Deprecated / no effect.} Originally intended as a grid-length control for \texttt{.skf}
files, but the present code ignores it when deciding whether an \texttt{.skf} is acceptable. Changing \is{ngrid}
does not influence the run.
\item[\is{grids}] Path to a grid file or directory copied into the rank-local scratch working directory before calling
\textbf{repopt} (\textbf{fit\_type = 0}). This is the only one of the three parameters that is directly used to stage
input for \textbf{repopt}.
\item[\is{rep.in}] Repopt input file. In \textbf{fit\_type = 0} this file is read and a rank-local
\texttt{rep.in} is created in scratch with optional additional blocks (\is{\$d3:}, \is{\$damph:}, \is{\$hubbardderivs:}, \is{\$damphver2:}).
\item[\is{libdir}] Output directory used to store generated \texttt{.skf} files. \textbf{erepopt} reads from
and writes to this directory during optimization.
\item[\is{libadddir}] Optional directory containing additional \texttt{.skf} files that are copied into the working
SKF directory for scoring (e.g. element pairs you do not want to optimize). If set, all \texttt{*.skf} in this
directory are copied for each evaluation.
\item[\is{repadddir}] Optional directory containing repulsive (Spline) tails appended to generated \texttt{.skf} files.
\item[\is{scratchfolder}] Base scratch directory. Each MPI rank uses subfolders
\texttt{slakos.tmp\_RANK}, \texttt{buildsks.tmp\_RANK}, and (for \textbf{fit\_type = 2}) \texttt{dftb.tmp\_RANK}.
\item[\is{skfclean}] If set to 1, \textbf{erepopt} deletes all files in \is{libdir} after each generation
(this reduces disk usage but can remove intermediate results).
\item[\is{outfile}] Main log file written by rank 0 (default: \texttt{gaserepfit.log}).
\item[\is{popinitialfile, popfinalfile}] Population snapshots written by rank 0. Each line is an individual:
all genes (radii + optional \is{\$d3:}/\is{\$damph:}/\is{\$hubbardderivs:}/\is{\$damphver2:}/\is{\$hubbards:}/\is{\$vorbes:})
followed by the individual score.
\end{description}

\textbf{How SKF files are reused and when repulsive tails are appended.}
During each evaluation, \textbf{erepopt} ensures that all required element-pair SKF files exist.
For an element pair \(X,Y\), it considers both directions (\texttt{X-Y.skf} and \texttt{Y-X.skf}).
The logic implemented in the current code is:

\begin{itemize}
\item \textbf{Prefer existing SKFs from \is{libadddir} (skip generation).}
If \is{libadddir} is set and a file \texttt{X-Y.skf} exists there, the pair is treated as available and
\textbf{skgen is not called} for that pair. (The same applies symmetrically for \texttt{Y-X.skf}.)
\item \textbf{Otherwise, check \is{libdir} for already-generated SKFs.}
If either \texttt{X-Y.skf} or \texttt{Y-X.skf} exists in \is{libdir} and passes an internal sanity check
(presence of required markers such as \texttt{Spline}, and when \is{omega} is enabled also matching \texttt{LC omega}),
the pair is treated as available and \textbf{skgen is not called}.
\item \textbf{Only if both directions are missing/invalid, call skgen once to generate both.}
When both \texttt{X-Y.skf} and \texttt{Y-X.skf} are missing (or fail the sanity check), \textbf{erepopt} runs
\texttt{skgen ... sktable X Y}. This single call produces both direction files, which are then moved into \is{libdir}.
\item \textbf{Repulsive tails from \is{repadddir} are appended only on the generation path.}
If \is{repadddir} is set and a matching tail file exists (e.g. \texttt{repadddir/X-Y.skf}), it is appended
to the newly generated SKF file(s) right after \texttt{skgen} finishes (in the rank-local scratch directory) and
the resulting SKF files are then moved into \is{libdir}. Therefore, \textbf{SKF files written to \is{libdir} by}
\textbf{erepopt already include the appended repulsive tails} when \is{repadddir} is enabled and the tail files exist.
If generation is skipped because an SKF
already exists in \is{libdir} or \is{libadddir}, \textbf{no additional repulsive-tail post-processing is performed}
on that existing file.
\end{itemize}

\textbf{Example:}
\begin{b4table}
  \$system:            &                          & &       \\
  \quad nthreads       &   1                      & &       \\
  \quad dftbversion    &   dftb+                  & &       \\
  \quad skgen          &   skgen                  & &       \\
  \quad onecent        &   hfatom\_spin           & &       \\
  \quad twocent        &   sktwocnt\_lr           & &       \\
  \quad repopt         &   repopt                 & &       \\
  \quad power          &   2                      & &       \\
  \quad dgrid          &   0.1                    & &       \\
  \quad ngrid          &   120                    & &       \\
  \quad grids          &   grids                  & &       \\
  \quad rep.in         &   rep4e.in               & &       \\
  \quad libdir         &   libskf4e               & &       \\
  \quad scratchfolder  &   /dev/shm               & &       \\
  \quad skfclean       &   0                      & &       \\
  \quad outfile        &   gaserepfit.log         & &       \\
  \quad popinitialfile &   pop.initial.dat        & &       \\
  \quad popfinalfile   &   pop.final.dat          & &       \\
  \$end                &                          & &       \\
\end{b4table}

\subsection{\$genetic\_algorithm:}

\begin{b4tableh}
  ga             & bool    &  0|1  &   1       \\
  runtest        & bool    &  0|1  &   0       \\
  fit\_type      & integer & 0|1|2 &   0       \\
  popsize        & integer &$\geq$1&   1000    \\
  preserved\_num & integer &$\geq$0&   100     \\
  ngen           & integer &$\geq$0&   1000    \\
  pmut           & integer &0.0:1.0&   0.02    \\
  pcross         & integer &0.0:1.0&   0.90    \\
  readr          & bool    &  0|1  &   1       \\
  restart        & bool    &  0|1  &   1       \\
\end{b4tableh}

\begin{description}
\item[\is{ga}] Switch GA on/off. In the current implementation, \textbf{erepopt} performs optimization
only when \is{ga}=1 and \is{runtest}=0.
\item[\is{runtest}] If 1, run in testing mode (no GA optimization).
\item[\is{fit\_type}] Select scoring mode:
0 = call \textbf{repopt} (repulsive fitting); 1 = eigenvalue fitting; 2 = DOS fitting.
\item[\is{popsize}] Population size.
\item[\is{preserved\_num}] Number of elite individuals preserved from generation \(n-1\) to \(n\).
\item[\is{ngen}] Number of generations.
\item[\is{pmut}] Mutation probability.
\item[\is{pcross}] Crossover probability.
\item[\is{readr}] If 1, the first initialization uses the values given in \is{\$element\_types:} and optional blocks
(\is{\$d3:}, \is{\$damph:}, etc.) as the initial genome; otherwise random initialization is used.
\item[\is{restart}] If 1, read \is{popfinalfile} and use its gene values as starting points.
\end{description}

\textbf{Example:}
\begin{b4table}
  \$genetic\_algorithm:     &            &&  \\
  \quad  ga                 &  1         &&  \\
  \quad  runtest            &  0         &&  \\
  \quad  fit\_type          &  0         &&  \\
  \quad  popsize            &  32        &&  \\
  \quad  preserved\_num     &  3         &&  \\
  \quad  ngen               &  30        &&  \\
  \quad  pmut               &  0.05      &&  \\
  \quad  pcross             &  0.9       &&  \\
  \quad  readr              &  1         &&  \\
  \quad  restart            &  0         &&  \\
  \$end:                    &            &&  \\
\end{b4table}


\subsection{\$variables: (optional)}

This block controls internal coupling rules used to constrain genes for \is{\$element\_types:} via the
\is{optype} field. These variables are mainly intended for advanced users and development workflows.

\begin{b4tableh}
  s1       & Real &  & 1.17 \\
  s2       & Real &  & 1.40 \\
  s3       & Real &  & 1.20 \\
  s4       & Real &  & 2.00 \\
  s5       & Real &  & 2.00 \\
  s6       & Real &  & 2.00 \\
  s7       & Real &  & 2.00 \\
  s8       & Real &  & 2.00 \\
  s9       & Real &  & 2.00 \\
  sdenswf  & Real &  & 2.00 \\
  idrefr   & integer & $\geq 0$ & 3 \\
\end{b4tableh}

\textbf{Note:} \is{idrefr} is an index into the flattened genome vector (0-based in code) used as a reference
gene by some \is{optype} coupling schemes.

\subsection{\$element\_types:}

This block defines the list of chemical elements to be optimized, together with their radii gene bounds,
initial values, and discretization precision. Each line has the form:

\begin{center}
  \texttt{ElementName  optype  lmax  (min r max prec)$_{lmax+2}$}
\end{center}

\begin{description}
\item[\is{ElementName}] Element symbol used by \texttt{skgen} (e.g. \texttt{H}, \texttt{C}, \texttt{O}). Case sensitive.
\item[\is{optype}] Integer controlling coupling constraints among radii genes (advanced usage; see \is{\$variables:}).
\item[\is{lmax}] Maximum angular momentum channel. The number of radii genes for this element is \(lmax+2\):
one density-compression radius (index 0) and \((lmax+1)\) wave-compression radii (S, P, D, ...).
\item[\is{min, r, max}] Minimum / initial / maximum value for each radius gene.
\item[\is{prec}] Precision level, used to set the mutation step size \(\delta\):
0 \(\rightarrow\) 1.0, 1 \(\rightarrow\) 0.1, 2 \(\rightarrow\) 0.01, 3 \(\rightarrow\) 0.001, 4 \(\rightarrow\) 0.0001, 5 \(\rightarrow\) 0.00001.
\end{description}

\textbf{Example:}
\begin{b9table}
  \$element\_types: & & &               &    &                &    &                 &    \\                 
  \quad H &  11 & 0     & 2.9  2.9  2.9 &  1 &  2.9  2.9  2.9 &  1 &                 &    \\                                                             
  \quad O & 111 & 1     & 2.7  2.8  2.9 &  1 &  2.7  2.8  2.9 &  1 &   2.7  2.8  2.9 &  1 \\
  \quad N & 111 & 1     & 3.0  3.2  3.4 &  1 &  3.0  3.2  3.4 &  1 &   3.0  3.2  3.4 &  1 \\
  \quad C & 111 & 1     & 3.6  3.8  4.0 &  1 &  3.6  3.8  4.0 &  1 &   3.6  3.8  4.0 &  1 \\
  \$end   &     &       &               &    &                &    &                 &    \\      
\end{b9table}

\subsection{\$atomparameters: (required for SK generation)}

This block provides per-element \texttt{AtomParameters} entries for \texttt{skdef.hsd}. Each element starts
with the element name on its own line, followed by a sub-block terminated by \textbf{\$subend}. Supported
keywords are:

\begin{description}
\item[\is{PowerCompression}] Integer power used in \texttt{PowerCompression}.
\item[\is{ShellResolved}] Value written to \texttt{ShellResolved}.
\end{description}

All other lines inside the element sub-block are copied verbatim into \texttt{skdef.hsd}.

\subsection{\$onecenterparameters: (required for SK generation)}

Lines inside this block are copied verbatim into the \texttt{OnecenterParameters} section of \texttt{skdef.hsd}.

\subsection{\$twocenterparameters: (required for SK generation)}

Lines inside this block are copied verbatim into the \texttt{TwoCenterParameters} section of \texttt{skdef.hsd}.

\subsection{\$addhamiltonian: (optional)}

Lines inside this block are appended verbatim inside the \texttt{Hamiltonian = DFTB \{ ... \}} section of the
rank-local \texttt{dftb\_in.hsd} file used for \textbf{fit\_type = 1/2}. This is the primary way to add DFTB+
Hamiltonian options (e.g. SCC parameters, dispersion, third-order terms, etc.).

\subsection{\$atomic\_energy: (optional; used in fit\_type=1)}

This block provides per-element atomic energy values used to compute a baseline energy \is{b0} for each compound
in \textbf{fit\_type = 1}. Each line is:
\begin{center}
  \texttt{element\_name  atomic\_energy}
\end{center}

\subsection{\$d3: (optional)}

This block defines D3-related scalar parameters optimized together with radii. Each line is:
\begin{center}
  \texttt{name  min  value  max  precision}
\end{center}

\textbf{Example:}
\begin{b4table}
  \$d3:    &                &    &  \\
  \quad s6 & 1.00 1.00 1.00 &  2 &  \\
  \quad s8 & 1.40 1.40 1.40 &  1 &  \\
  \quad a1 & 0.48 0.48 0.48 &  2 &  \\
  \quad a2 & 4.70 4.70 4.70 &  1 &  \\
  \$end:   &                &    &  \\
\end{b4table}

\subsection{\$damph: (optional)}

Damping parameter optimized as a single scalar. The line format is:
\begin{center}
  \texttt{min  value  max  precision}
\end{center}

\subsection{\$hubbardderivs: and \$hubbardderivsfull: (optional)}

Hubbard derivative parameters optimized as per-element scalars. Each line is:
\begin{center}
  \texttt{name  min  value  max  precision}
\end{center}

\textbf{Note:} \is{\$hubbardderivsfull:} enables a different internal flag but uses the same line format.

\subsection{\$damphver2: (optional)}

Second version of damping parameters, optimized as per-element scalars. Each line is:
\begin{center}
  \texttt{name  min  value  max  precision}
\end{center}

\subsection{\$hubbards: (optional)}

Customized Hubbard parameters per element. Each line is:
\begin{center}
  \texttt{name  type  min  value  max  precision}
\end{center}

\subsection{\$vorbes: (optional)}

Customized on-site parameters per element. Each line is:
\begin{center}
  \texttt{name  type  min  value  max  precision}
\end{center}
  
\textbf{Example:}
\begin{b4table}
  \$vorbes:&    &                     &     \\
  \quad  N & 2S & -0.83  -0.82  -0.81 &  3  \\
  \$end:   &    &                     &     \\
\end{b4table}

\subsection{\$compounds: (required for fit\_type=1/2)}

This block defines the training set. Each line is:
\begin{center}
  \texttt{xyzfile  charge  start  end  gauss  reference\_data  weight}
\end{center}

\begin{description}
\item[\is{xyzfile}] Geometry in XYZ format. Element names in the XYZ are lowercased internally.
\item[\is{charge}] Total charge passed to DFTB+.
\item[\is{start, end}] Energy window (in the same units as the reference data) used for scoring.
\item[\is{gauss}] Gaussian broadening used by \texttt{dp\_dos} (fit\_type=2) and stored for logging.
\item[\is{reference\_data}] Reference file:
in \textbf{fit\_type = 1} it is read as a list of eigenvalues and occupations;
in \textbf{fit\_type = 2} it is read as a DOS curve (energy, DOS).
\item[\is{weight}] Weighting factor. For \textbf{fit\_type = 2}, special negative values implement
energy-dependent weighting: \(-1\) linear, \(-2\) quadratic, \(-4\) quartic weighting from start to end.
\end{description}

\section{Output}

\textbf{erepopt} writes several outputs during a GA run:

\begin{description}
\item[\is{Main log (\is{outfile})}] Written by rank 0. It contains progress lines during evolution and ends with
``\textbf{** erepobj normal termination **}'' on success. After GA completion it prints an ``interpreted input''
summary, the final optimized parameters, and per-compound errors.

\item[\is{score.dat}] GA score history written by GALib (frequency controlled by \is{scoref}/\is{flushf}).
\item[\is{\is{popinitialfile}, \is{popfinalfile}}] Population snapshots written by rank 0 at initialization and
after every generation.
\item[\is{Generated SKF files}] Stored in \is{libdir}. File names encode the element radii and optional parameter
blocks, e.g. \texttt{X\_d...\_w...-Y\_d...\_w...\_hub...\_vor....skf} (and an additional \texttt{omega\_...} prefix if \is{omega} is set).
\item[\is{Temporary scratch}] Rank-local folders under \is{scratchfolder} are created and deleted during the run.
\end{description}

\section{Tips}

\begin{description}
\item[\is{Run under a scheduler}] Since \textbf{erepopt} uses MPI, on clusters you typically run it with
\texttt{srun} (recommended) or \texttt{mpirun}. Example (2 MPI ranks, 8 threads each):
{\scriptsize
\begin{verbatim}
srun -n 2 -c 8 erepopt erep.inp
\end{verbatim}}

\item[\is{Validate tool paths first}] Ensure \is{skgen}, \is{onecent}, \is{twocent}, \is{repopt}, \is{dftbversion},
and (for \textbf{fit\_type = 2}) \texttt{dp\_dos} are available in \texttt{\$PATH} or specified with absolute paths.

\item[\is{Start from a small training set}] Debug input formats and external tool calls with 1--2 compounds and a small
\is{popsize}/\is{ngen} before scaling up.

\item[\is{Be careful with skfclean}] Setting \is{skfclean}=1 removes all \texttt{.skf} files in \is{libdir} each generation,
which can make debugging difficult. Consider using \is{skfclean}=0 while developing.

\item[\is{Common failure modes}] Most runtime errors come from missing files or failing external calls.
If you see an error mentioning \texttt{rep.out} or missing DOS output, inspect the corresponding scratch directory
for the rank and check the external program logs (\texttt{rep.out}, \texttt{log}, \texttt{dos.dat}).
\end{description}


